\scanpage{136}
\mspage{81}
% archive sheet number: 135

comme \inserted{\unsure{(,)}} provenant de $X,Y$ de $\mathrm{Cat}_{/S}$:

\[
\boxed{\mathrm{Hom}_{W_S^d}(\deleted{\unsure{C}}X,CY)\simeq\mathrm{Hom}_{W_S^g}(X,CY)\simeq\mathrm{Hom}_{W_S}(X,Y),}
\]

% slanted
\margin{Cette formule suggère que si $X\in\mathrm{Ob}\,\mathrm{Cat}_{/S}$, $Y\in\mathrm{Ob}\,\underline{\mathrm{Fibhot}}\,S$, alors $\mathrm{Hom}_{W_S^d}(X,Y)\unsure{\simeq_{?}|\simeq}\mathrm{Hom}_{W_S^g}(X,Y)\simeq\mathrm{Hom}_{W_S}(X,Y)$, mais ça résulte de la formule d'adjonction et de $Y\to C(Y)$ dans $W_S^d\cap W_S^g$ pour $Y$ dans $\underline{\mathrm{Fibhot}}\,S$.}

où $\mathrm{Hom}_{W_S^d}$, $\mathrm{Hom}_{W_S^g}$ et $\mathrm{Hom}_{W_S}$ désignent respectivement les \deleted{\unsure{ens.}} $\mathrm{Hom}$ dans les catégories localisées par $W_S^d$, $W_S^g$ et $W_S$, i.e. dans $\mathrm{HOT}_W(S)$, $\mathrm{HOT}_W^{\unsure{\circ|o}}(S)$ et $\mathrm{HOT}_W^{\unsure{\mathrm{lc}|\mathrm{le}}}(S)$ respectivement.

\uline{Commentaire} Toutes les constructions et les énoncés de \textcircled{I} à \textcircled{VI}, à l'exception du fait que les foncteurs de descente induisent des équiv. pour les localisés \uline{grossiers} (pour \deleted{\unsure{les $S$}} les localisés grossiers), \inserted{et des formules d'adjonction de \textcircled{IV},} sont de nature quasiment soritale, et n'exigent d'autre \deleted{\unsure{propr.}} condition sur le localiseur \unsure{fondt|fondtl} $\unsure{W|\mathcal{W}}\subset\mathrm{Fl}\,\mathrm{Cat}$ que \deleted{\unsure{les celles de}} L1, L3 \unsure{faible}, L4 bis, récapitulées comme a) b) c) de la page 45. (On n'a même pas à utiliser d).) Il est vrai qu'en l'absence de la condition L5 (page 43), on ignore si à l'étage les deux localiseurs pertinents, $W_S$ et $W_S^f$ (égal à $W_S^d$ du côté gauche, à $W_S^g$ du côté droit) sont égaux. Donc à priori on trouve à l'étage \uline{deux} catégories localisées \unsure{évt|cert.|a priori} distinctes, et non pas une seule, savoir $(\underline{\mathrm{Fibhot}}\,S)\Sigma^{-1}$ avec $\Sigma=W_S$ ou $\Sigma=W_S^f$ $(=W_S^d=W_S^g)$.

Par contre, le fait que les foncteurs de descente induisent des équivalences entre les localisés grossiers, est plus profond, il exige \deleted{\unsure{la f}} \inserted{(me semble-t-il)} qu'on connaisse un foncteur ``montant'' $C$, disons de $\mathrm{Cat}_{/S}$ vers $\underline{\mathrm{Fibhot}}\,S$, \deleted{\unsure{ayant transformant les}} compatible avec $W_S$, et muni de

\[
X\xrightarrow{i_X}C(X)
\]

qui soit de $W_S$ pour $X$ quelconque. Quand on sait que $\mathrm{Cat}$ est une cat. de Quillen, même sans avoir les propriétés du foncteur $\underline{\mathrm{Ch}}_\infty$ ou $\underline{\mathrm{Ch}}$,\textsuperscript{*} on doit sans doute

% slanted
\margin{* \unsure{ni|ou} même être sûr de l'existence d'une factorisation \uline{fonctorielle} de $X\xrightarrow{f}S$ en $X\xrightarrow{i_X}C(X)\xrightarrow{\bar f}S$, \unsure{où} $i_X\in W$ et $\bar f$ \unsure{parf.}}
